\section{Related Work}
\label{sec:related}

Our work sits at the intersection of three literatures: sycophancy measurement, persona assignment in LLMs, and recent friction-as-alignment proposals. We summarize each and position our contribution.

\para{Sycophancy in LLMs.} \citet{perez2022discovering} first established that RLHF makes models more sycophantic on stated-belief questions, with the effect growing with scale. \citet{sharma2023sycophancy} extended the analysis to five production assistants (Claude 1.3/2, GPT-3.5, GPT-4, LLaMA-2-70b-chat) across four behavior types -- feedback, ``are you sure?'', answer, and mimicry sycophancy -- and released \sycoeval. They also showed by Bayesian logistic regression on hh-rlhf that ``matches user's beliefs'' is the single most predictive feature of human preference judgments, implicating preference-data design as the structural cause. \citet{wei2023simple} proposed a synthetic-data fine-tuning fix that reduces sycophancy by 8--22\% on three NLP tasks. \citet{liu2025truthdecay} extend the analysis to multi-turn dialog and show that sycophancy compounds with conversational pressure across 1, 3, and 7 turns. Unlike these works, we hold the model fixed and ask whether \emph{persona} -- a one-line change to the system prompt -- can reduce the failure rate.

\para{Persona assignment and bias.} \citet{deshpande2023toxicity} showed that persona prompts can amplify toxicity up to six-fold, an early demonstration that persona is not cosmetic. \citet{gupta2023bias} ran the largest scale study (19 personas \(\times\) 24 reasoning datasets \(\times\) 4 LLMs) and found 80\% of personas cause significant accuracy drops on at least one task, with worst-case relative drops above 70\%. Persona-induced bias is a primary obstacle to using persona as an alignment tool. \citet{wang2023rolellm} curated \textsc{RoleBench} (168k samples, 100 roles) for evaluating persona \emph{fidelity}; \citet{zhang2018personachat} provide an older dialog-grounded persona dataset. Our work treats persona orthogonally: rather than test whether persona harms reasoning, we test whether a particular kind of persona -- one specifying epistemic commitments -- helps the model resist user pressure.

\para{The Personality Illusion.} The most direct counter-evidence to our hypothesis is \citet{han2025personality}, who run three RQs across twelve open-source LLMs and find that instructional alignment stabilizes Big-Five self-reports (variability drops 40--45\%) and that persona injection moves these self-reports as predicted (\(\beta=3.95\) for agreeableness). But persona injection does not move behavior: only 24\% of trait-task associations are significant and a sycophancy-targeting persona moves sycophantic behavior by \(\beta=0.03\) (\(p=0.67\)). Their conclusion -- that current personas are linguistic surface, not action-controlling identity -- is the strong claim our work tests. We replicate the lack of behavioral effect for demographic personas and refute it for belief-based personas.

\para{Constitutional AI and frictive alignment.} \citet{bai2022constitutional} introduce Constitutional AI: instead of human harm labels, the model critiques and revises its outputs against a natural-language constitution. Constitutional models can produce principled refusals, but \citet{sharma2023sycophancy} show they still fold under user-stated belief pressure. \citet{huang2024pinkelephants} (Direct Principle Feedback) is a lighter-weight CAI alternative trained via DPO on critique-revise pairs. \citet{pustejovsky2026frictive} propose \emph{frictive policy optimization}, an alignment framework treating clarification, verification, challenge, redirection, and refusal as first-class control actions; their approach is theoretical and lacks an empirical instantiation we could find. Our constitutional condition is a faithful prompt-only analog of CAI, and our belief condition operationalizes one variant of the friction-as-stance ideas of \citet{pustejovsky2026frictive}.

\para{Role-play, simulacra, and disagreement.} \citet{shanahan2023roleplay} argue that an LLM dialogue agent is best understood as role-playing a character, or as a superposition of simulacra; the low-identity default is one point in role-space. \citet{du2023debate} show that debating LLM instances improves factuality, but a recent controlled follow-up (``Can LLM Agents Really Debate?'', 2025) shows much of this gain is ensembling rather than genuine disagreement: forced disagreement degrades performance. This is indirect evidence for our hypothesis -- without distinct persona, two instances of the same base model converge. \citet{sun2023delphi} provide \textsc{Delphi}, a controversial-topic benchmark relevant for measuring stance-taking; \citet{zhou2023sotopia} provide \textsc{Sotopia}, a goal-directed social-interaction benchmark.

\para{Position of this work.} No published study has tested whether substantive belief-based personas (as opposed to demographic, role-based, or Big-Five) close the dissociation between persona self-reports and persona behavior documented by \citet{han2025personality}. Doing so is our central contribution.
